\maketitle%与上面的\title对应
\thispagestyle{plain}%摘要专用页显示页码 1，并与后续正文连续编号
\begin{abstract}%摘要+摘要+摘要+摘要+摘要

	本文以\textbf{信息集边界}为主线，将四个微网调控问题统一到母线能量平衡、储能状态递推、容量／功率边界与弃光约束构成的公共模型中；各问仅改变可用信息与结算规则。每次滚动决策都继承上一执行块的实测储电量，未启用节点则延续最近有效计划，从而严格保持时序一致性。

	\textbf{问题一}为确定性日前调度。以全天购电费最小为目标、首末储电量均为 $6000$ kWh 为闭环约束建立线性规划，并用“费用—吞吐量—弃光量”三层词典序目标消除退化最优解。最优全天购电量 \QoneEnergy\ kWh，购电费 \QoneCost\ 元，较无储能基准的 \QoneNoStorage\ 元节省 \QoneSave\ 元（\QoneSavePct\%）。

	\textbf{问题二}采用只依赖历史观测的负载—光伏因果预测和逐日确定性等价优化。$5$ 倍紧急电价在孤立单时段下对应 $4:1$ 非对称损失及光伏输出 $0.2$ 分位启发；全年模型通过样本外费用在多个残差分位候选中在线选择，并用“先撤充电、再增放电、最后紧急购电”的实时反馈执行。$2$ 月至 $12$ 月总费用为 \QtwoTotal\ 元，紧急购电 \QtwoEmergency\ kWh。对照实验表明，把当天真实负载误作已知会低估 \DeltaLoad\ 元（\DeltaLoadPct\%）。

	\textbf{问题三}按题意采用“计划费 $+$ 调整费 $+$ 紧急费”的叠加结算，并用 $0$--$1$ 互斥约束禁止同时充放电。仅 $0{:}00$ 决策时费用为 \QthreeSzero\ 元，依次加入 $6{:}00$、$12{:}00$、$18{:}00$ 更新后降至 \QthreeSone、\QthreeStwo、\QthreeSthree\ 元，三节点均有正收益，累计节省 \QthreeVtotal\ 元。配对控制把其中 \QthreeLoadTotal\ 元归因于负载观测刷新、\QthreePvTotal\ 元归因于新光伏预报，避免把综合节点收益误称为纯光伏价值。

	\textbf{问题四}按题面直接代入附件 4 电价重算，Q4-2 与 Q4-3 主结果分别为 \QfourTwoMain\ 元和 \QfourThreeMain\ 元，滚动更新节省 \QfourRollPerfect\ 元。另设因果价格扩展，费用分别为 \QfourTwoCausal\ 元和 \QfourThreeCausal\ 元；两种价格信息情形下滚动更新均有效。

	全部指定日期表均由最终结果工作簿自动生成。物理残差、信息泄漏 mutation、结算分解、计划持续性与回归测试全部通过；主方案同时充放电量与费用分解回代差均低于 $10^{-8}$。

	\keywords{微网经济调度；因果预测；滚动优化；混合整数规划；信息集}
\end{abstract}
